\documentclass[12pt,a4paper]{report}

\usepackage{siunitx}
\usepackage{tcolorbox}
\usepackage[utf8]{inputenc}
\usepackage[T1]{fontenc}
\usepackage{tikz}
\usepackage{amsmath, amssymb, amsthm}
\usepackage{graphicx}
\usepackage{hyperref}
\usepackage{geometry}
\usepackage{booktabs}
\usepackage{listings}
\usepackage{xcolor}
\usepackage{float}
\usepackage{titlesec}
\usepackage{fancyhdr}
\usepackage{enumitem}

\setlength{\parindent}{0pt}

\renewcommand{\thesection}{\arabic{section}}

\usetikzlibrary{calc}

% Page Setup
\geometry{margin=1in}
\hypersetup{colorlinks=true, linkcolor=blue, urlcolor=cyan}

% Header Footer
\pagestyle{fancy}
\fancyhf{}

\lhead{\includegraphics[height=0.8cm]{Image/uni_logo_old.jpg}}

\rhead{
\begin{tabular}{r}
University of Science \\
Faculty of Information Technology
\end{tabular}
}

\cfoot{\thepage}

\renewcommand{\headrulewidth}{0.4pt}

% Code Snippet Styling
\definecolor{codegreen}{rgb}{0,0.6,0}
\definecolor{codegray}{rgb}{0.5,0.5,0.5}
\definecolor{codepurple}{rgb}{0.58,0,0.82}

\lstset{
language=C++,
basicstyle=\ttfamily\small,
keywordstyle=\color{blue},
stringstyle=\color{codepurple},
commentstyle=\color{codegreen},
breaklines=true,
numbers=left,
numberstyle=\tiny\color{codegray},
frame=single
}

\setlength{\parindent}{0pt}

\begin{document}

% =====================================================
% TITLE PAGE
% =====================================================

\begin{titlepage}

\begin{tikzpicture}[overlay,remember picture]

\draw [line width=2.25pt]
($ (current page.north west) + (1,-1) $)
rectangle
($ (current page.south east) + (-1,1) $);

\draw [line width=1pt]
($ (current page.north west) + (1.2cm,-1.2cm) $)
rectangle
($ (current page.south east) + (-1.2cm,1.2cm) $);

\end{tikzpicture}

\begin{center}

\textbf{\large VNUHCM - UNIVERSITY OF SCIENCE}

\vspace{0.2cm}

\textbf{\large FACULTY OF INFORMATION TECHNOLOGY}

\vspace{0.8cm}

\includegraphics[width=0.3\textwidth]{Image/uni_logo_old.jpg}

\vspace{0.8cm}

{\fontsize{36}{42}\selectfont \textbf{Research Documentation}}

\vspace{0.6cm}

{\Huge Note Book}

\vspace{2cm}

{\Large \textbf{General Information}}

\vspace{0.6cm}

\begin{tabular}{rl}

Name: & Nguyen Phu Trong \\

Student ID: & 25125067 \\

\end{tabular}

\vfill

{\large \today}

\end{center}

\end{titlepage}

\tableofcontents 

\newpage % ================= PART ================= 

\setcounter{part}{0}
% Đảm bảo khai báo các gói này ở phần Preamble (trước \begin{document}):
% \usepackage{amsmath, amssymb}
% \usepackage{graphicx}
% \usepackage{float}
% \usepackage{enumitem}

\chapter{Machine Learning}

\section{Introduction}
So, what is \textbf{Machine Learning}?

\begin{quote}
    \emph{Machine Learning is the field of study that gives computers the ability to learn without being explicitly programmed.}
\end{quote}

Based on the nature of the learning tasks, Machine Learning can be broadly categorized into three main paradigms:
\begin{itemize}[itemsep=2pt, topsep=4pt]
    \item \textbf{Supervised Learning}
    \item \textbf{Unsupervised Learning}
    \item \textbf{Reinforcement Learning}
\end{itemize}

\subsection{Supervised Learning}

\paragraph{Example Problem: House Price Prediction}
\begin{description}[style=nextline, leftmargin=1.5cm, itemsep=4pt]
    \item[Given:] A dataset $\mathcal{D}$ containing $n$ labeled samples:
    \begin{equation*}
        \mathcal{D} = \left\{ \left(x^{(1)}, y^{(1)}\right), \left(x^{(2)}, y^{(2)}\right), \dots, \left(x^{(n)}, y^{(n)}\right) \right\}
    \end{equation*}
    \item[Task:] Given a new residence with an area of $x$ square feet, predict its corresponding price $\hat{y}$.
\end{description}

\begin{figure}[H]
    \centering
    \includegraphics[width=0.7\linewidth]{Image/HousePricePrediction.png}
    \caption{Illustration of the house price prediction problem.}
    \label{fig:house_price_prediction}
\end{figure}

\subsubsection*{Regression vs. Classification}
Depending on the nature of the target variable $y$, supervised learning problems are further divided into two main types:

\begin{description}[leftmargin=0.5cm, itemsep=6pt]
    \item[Regression:] The target variable $y \in \mathbb{R}$ is \textbf{continuous}. 
    \newline \emph{Example:} Predicting house prices, where the output can take any real measurable value (e.g., \$2.5 billion or \$3.6 billion).
    
    \item[Classification:] The target label $y$ is \textbf{discrete} (categorical), taking values from a finite set $y \in \{1, 2, \dots, C\}$. 
    \newline \emph{Example:} Predicting whether an incoming email is spam ($y = 1$) or non-spam ($y = 0$).
\end{description}

\begin{figure}[H]
    \centering
    \includegraphics[width=0.7\linewidth]{Image/ClassificationProblem.png}
    \caption{An example for classification problem.}
    \label{fig:placeholder}
\end{figure}

\subsubsection*{High-dimensional Inputs}

\begin{align*}
    x \in \mathbb{R^{$d$}} \text{ for large $d$}
\end{align*}

E.g., 

\begin{align*}
    x =\begin{bmatrix} 
        x_1 \\ 
        x_2 \\
        x_3 \\ 
        ... \\
        ... \\
        x_d
    \end{bmatrix} \rightarrow y: \text{price}
\end{align*}

Where, $x_1, x_2, x_3, ..., x_d$ are the properties (living size, lot size, # floors, condition, ...)

\subsubsection{Supervised Learning in Computer Vision}

\begin{description}[leftmargin=0.5cm, itemsep=6pt]
    \item [Image Classification] x = raw pixels of the image (input), y = the main object (output) 
\end{description}

\begin{figure}[H]
    \centering
    \includegraphics[width=0.5\linewidth]{Image/ImageClassification.png}
    \caption{ImageNet Large Scale Visual Recognition Challenge. Russkovshy et al.'2015}
    \label{fig:placeholder}
\end{figure}

\subsubsection{Supervised Learning with Neural Nets(Deep Learning)}

\begin{figure}[H]
    \centering
    \includegraphics[width=0.8\linewidth]{Image/NeuralNets.png}
    \label{fig:placeholder}
\end{figure}

\subsection{Unsupervised Learning}

\begin{description}[style=nextline, leftmargin=1.5cm, itemsep=4pt]
    \item[Given: ] Dataset contains \textbf{no labels}: $x^{(1)}, x^{(2)}, ..., x^{(n)}$ 
    \item[Details: ] These data set just included input and there is no output, you just have a bunch of data and you don't know what the type of the output is. 
\end{description}
\begin{figure}[H]
    \centering
    \includegraphics[width=0.8\linewidth]{Image/Supervised_And_ItsCounterpart.png}
    \caption{Different between supervised and unsupervised through House Price Predict}
    \label{fig:placeholder}
\end{figure}
\subsubsection{Large Language Models}

\begin{description}
    \item[Pretraining: ] General-purpose models learnt on unlabled, massive datasets. 
    - transformers architecture
    - next-word prediction loss
    Prompting, instruction tuning, think step by step
\end{description}
\subsection{Reinforcement Learning}

\begin{description}
    \item[Purpose: ]
        Learning to make  sequential decisions.
        Train stochastic models with rewards.
    \item[Details: ]
        The algorithm can collect data interactively
\end{description}
\begin{figure}
    \centering[H]
    \includegraphics[width=0.8\linewidth]{Image/ReinforcementLearning.png}
    \label{fig:placeholder}
\end{figure}

\subsubsection{RL for Training LLMs}

\begin{figure}[H]
    \centering
    \includegraphics[width=0.5\linewidth]{Image/RLforLLMs.png}
    \label{fig:placeholder}
\end{figure}

\begin{itemize}
    \item Use the reward as a guidance for updating the LLM to generate better texts
    \item RL is needed partly because the LLM generation are stochastic and thus not directly differentiable 
\end{itemize}

\section{Supervised Learning Setup}

A \textbf{hypothesis} or a prediction function is function: $h:$ $X$ $\rightarrow $ $Y$

\begin{itemize}
    \item $X$ is an image, $Y$ contains "cat" or "dog".
    \item $X$ is a text snipper, $Y$ contains "hate speech" or "not".
\end{itemize}

A \textbf{training set} is a set of pair:
\[
    (x^{(1)}, y^{(1)}), (x^{(2)}, y^{(2)}), ... , (x^{(n)}, y^{(n)})
\]

s.t. $x^{(i)} \in X$ and $y^{(i)} \in Y$ for i = 1, ..., n

\begin{description}
    \item[Given: ] A training set our goal is to produce a \texttt{good} prediction function $h$ 
    
    - Defining "good" will take us a bit. It's a modeling question!
    
    - We will want to use $h$ on new data not in the training set
\end{description}


We can define $h$ as a linear function:
\[
    h = \theta_0 + \theta_1x
\]

\begin{figure}[H]
    \centering
    \includegraphics[width=0.8\linewidth]{Image/LinearFitonHousingData.png}
    \label{fig:placeholder}
\end{figure}


So that, we are trying to minimize the different between these points to the line.

Now, we add features (bedrooms and lot size) to incorporate more information about houses:

\begin{figure}[H]
    \centering
    \includegraphics[width=0.8\linewidth]{Image/boom.png}
    \label{fig:placeholder}
\end{figure}

What is the prediction?
\[
    h(x) = \theta_0 + \theta_1x_1 + \theta_2 x_2 + \theta_3 x_3
\]

With the convention that $x_0 = 1$, we can rewrite the expression as:
\[
    h(x) = \sum^{3}_{i = 0} \theta_i x_i
\]

We write the vectors as (important notation):
\begin{align*}
    \theta = \begin{bmatrix}
        \theta_0 \\ \theta_1 \\ \theta_2 \\ \theta_3 
    \end{bmatrix} \text{ and } 
    x^{(1)} = \begin{bmatrix}
        x_0^{(1)} \\ x_1^{(1)} \\ x_2^{(1)} \\ x_3^{(1)} 
    \end{bmatrix} 
            = \begin{bmatrix}
            1 \\ 2104 \\ 4 \\ 45
    \end{bmatrix} \text{ and } y^{(1)} = 400
\end{align*}


We call $\theta$ \textbf{parameters}, $x^{(i)}$ is the input or the \textbf{features}, and the output or \textbf{target} $y^{(i)}$. To be clear:

$(x, y)$ is a training example, and $(x^{(i)}, y^{(i)})$ is the $i^{th}$ example.

\begin{description}
    \item[Least Square: ]
        
        Let $h_\theta(x) = \sum^{d}_{j=0} \theta_jx_j$ want to choose $\theta$ so that $h_0 \approx y$.

        The idea is: $J(\theta) = \frac{1}{2}\sum^{n}_{i = 1} (h_\theta(x^{(i)}) - y^{(i)})^2$

    \item[Solving Least Square Optimization Problems: ] 
        
        Using Stochastic Gradient Descent for simple data:
        
        Define: \begin{align*}
            h(x) &= \theta_0 + \theta_1 x \\
            e &= \frac{1}{2}(\hat{y} - y)^2 
        \end{align*}
        
        Then, \begin{equation*}
            \frac{\partial e}{\partial\theta_0} = \frac{\partial e}{\partial y} \cdot \frac{\partial y}{\partial \theta_0} = (\hat{y}  - y) \cdot 1 = \hat{y} - y
        \end{equation*}

        \begin{equation*}
            \frac{\partial e}{\partial\theta_1} = \frac{\partial e}{\partial y} \cdot \frac{\partial y}{\partial \theta_1} = (\hat{y}  - y) \cdot x_j
        \end{equation*}
        
        \begin{figure}[H]
            \centering
            \includegraphics[width=0.8\linewidth]{Image/StochasticGradientDescent.cpp.png}
            \caption{Details implementation}
            \label{fig:placeholder}
        \end{figure}
\end{description}

\section{Introduction to Classification}

\subsection{Logistic Regression: Link Functions}

\begin{description}
    \item[Given :] A training set $\{x^{(i)}, y^{(i)}\}$ for $i = 1. ..., n$  let $y^{(i)} \in \{0, 1\}$. Want $h_\theta(x) \in [0, 1]$. Let's pick a smooth function: 
    \[
        h_\theta(x) = g(\theta^T x)
    \]
    Here $g$ is a link function. There are many ... but we'll pick one. 
    \[
        g(z) = \frac{1}{1 + e^{-z}}
    \]
    \item [How do we interpret $h(x)$ ?]:
    
        $P(y = 1 | x, \theta) = h_\theta(x)$ 

        $P(y = 0 | x, \theta) = 1 - h_\theta(x)$

    Let's write the Likelihood function: 
    \[
        L(\theta) = P(y | X; \theta) = \prod^{n}_{i = 1} p(y^{(i)} | x^{(i)}; \theta) = \prod^{n}_{i = 1} h_\theta(x)^{y^{(i)}} (1 - h_\theta(x))^{1 - y^{(i)}}
    \]

\end{description}

\subsection{Newton's Method}

\subsection{Summary}

\begin{figure}[H]
    \centering
    \includegraphics[width=0.8\linewidth]{Image/OptimizationMethodSummary.png}
    \label{fig:placeholder}
\end{figure}
\textbf{Why Mini-batch SGD Is the Practical Standard}

\begin{itemize}
    \item \textbf{Hardware Optimization and Parallelism:} 
    Modern linear algebra libraries and hardware accelerators (GPUs, CPUs with SIMD/AVX instruction sets) are highly optimized for matrix operations. Processing a batch of size $B \times D$ achieves superior parallel throughput, with execution latency comparable to processing an individual sample ($1 \times D$) sequentially.

    \item \textbf{Faster Convergence Rate Than Batch Gradient Descent:} 
    The model does not need to compute gradients across all $N$ samples in the dataset before executing a parameter update. Instead, parameters $\theta$ are updated iteratively after evaluating each small mini-batch of size $B$, substantially reducing wall-clock training time on large-scale datasets.

    \item \textbf{Balance Between Gradient Stability and Generalization:} 
    The batch-averaged gradient vector $g = \frac{1}{B} \sum_{i=1}^B \nabla_\theta \mathcal{L}(h_\theta(x^{(i)}), y^{(i)})$ significantly dampens the high variance and noise inherent in pure SGD ($B = 1$), producing a smoother optimization path. Simultaneously, the stochastic fluctuations across successive mini-batches inject sufficient variance to escape saddle points and suboptimal local minima.
\end{itemize}

\section{Exponential Family \& GLMs Classification}

This section focuses on proofs; therefore, I will refer back to it whenever I need further clarification or deeper insight into this field.

\section{Generative Learning Algorithms}


\section{Bias, Variance, Regularization, and Model Selection}

\subsection{Fundamental Tension - Overfit and Underfit}

\textbf{A central question in ML: } given a finite, noisy sample, how do we pick a model that will generalize?

\begin{description}
    \item[The setup: ] We never see the True Function

    There is a true function (say $h_\theta(x) = \theta_2x^2 + \theta_1 x + \theta_0$) that generates the world. We never see it. We see a finite set of noisy samples:
    \[
        y^{(i)} = h_\theta(x^{(i)}) + \varepsilon^{(i)}
    \]
    
\end{description}

\begin{figure}[H]
    \centering
    \includegraphics[width=0.8\linewidth]{Image/setup.png}
    \label{fig:placeholder}
\end{figure}

\begin{tcolorbox}[
    colback=gray!5,
    colframe=gray!70!black,
    arc=3mm,
    boxrule=0.8pt,
    left=5mm,
    right=5mm,
    top=3mm,
    bottom=3mm
]
Every fresh dataset $S$ is a different set of points.
We should judge our model across all the datasets
we could have seen.
\end{tcolorbox}

\begin{description}
    \item[Three Different Samples $S^{(1)}, S^{(2)}, S^{(3)}$: ] Same true function. Three independent draws of the training data.  
\end{description}

\begin{figure}[H]
    \centering
    \includegraphics[width=0.8\linewidth]{Image/ThreeDifferentSamples.png}
    % \caption{Caption}
    \label{fig:placeholder}
\end{figure}

We will try to fit the same model class to each one and see how much the fit changes.

\begin{description}
    \item[Fit a line: ] Too simple 
    
    \begin{figure}[H]
        \centering
        \includegraphics[width=0.8\linewidth]{Image/aline.png}
        \label{fig:placeholder}
    \end{figure}

    \item [High bias/ underfit. ] Fix: use a more expressive model class.
\end{description}


\begin{description}
    \item[Fit a High Degree Polynomial: ] Too flexible

    \begin{figure}[H]
        \centering
        \includegraphics[width=0.8\linewidth]{Image/highpoly.png}
        \label{fig:placeholder}
    \end{figure}

    \item [High bias / ovefit .] Fix: use a less expressive class with lower variance
\end{description}

\begin{description}
    \item[Fit a Quadratic: ] Just right

    \begin{figure}[H]
        \centering
        \includegraphics[width=0.8\linewidth]{Image/quadratic.png}
        % \caption{Caption}
        \label{fig:placeholder}
    \end{figure}

    \item [The fits: ] similar across samples and stay close to the data.

    \item [Takeaway: ] It fits! The "right" model class balances expressiveness against sensitivity to the particular sample we drew
\end{description}

\subsection{Training vs. Test Error as We Change Complexity}

\begin{figure}[H]
    \centering
    \includegraphics[width=0.8\linewidth]{Image/error.png}
    % \caption{Caption}
    \label{fig:placeholder}
\end{figure}

Test error = $\text{bias}^2$ + variance + error. Training error sits below.

\begin{tcolorbox}[
    colback=gray!5,
    colframe=gray!70!black,
    arc=3mm,
    boxrule=0.8pt,
    left=5mm,
    right=5mm,
    top=3mm,
    bottom=3mm
]
Model selection must minimize \textbf{expected test error}, not training error.
\end{tcolorbox}

\subsection{Same story in classification}

This is not regression-only phenomenon. In classification, an overfit decision surface is a wiggly boundary that  perfectly separates the training points but generalizes poorly

\begin{figure}[H]
    \centering
    \includegraphics[width=0.8\linewidth]{Image/classification.png}
    \caption{Caption}
    \label{fig:placeholder}
\end{figure}

\subsection{The Bias - Variance Decomposition}

Bias--Variance Decomposition helps us understand
the sources of prediction error by separating
the expected squared error into three components:

\[
\mathbb{E}[(Y-\hat{f}(x))^2]
=
\underbrace{\operatorname{Bias}^2}_{\text{Systematic Error}}
+
\underbrace{\operatorname{Variance}}_{\text{Model Sensitivity}}
+
\underbrace{\sigma^2}_{\text{Irreducible Noise}}.
\]

It explains why a model may fail to generalize:
high bias leads to underfitting, while high variance
is associated with overfitting.

The goal is to understand the tradeoff between
bias and variance and minimize the expected
prediction error on unseen data.

\subsection{From Bias--Variance Decomposition to Fixes}

\subsubsection{Regularization}

Regularization refers to techniques used to reduce
overfitting and improve generalization by constraining
the effective complexity of a model.

A common approach is to add a penalty term to
the training loss:

\[
\mathcal{L}_{\text{total}}
=
\mathcal{L}_{\text{train}}
+
\lambda \Omega(\theta),
\]

where $\Omega(\theta)$ is a regularization penalty
and $\lambda$ controls its strength.

Regularization typically reduces variance at the
cost of introducing some additional bias.

\subsubsection{From Decomposition to Fixes}

The bias--variance decomposition motivates two
important techniques:

\begin{enumerate}
    \item \textbf{Dev / Holdout Set:}
    A separate validation dataset provides an
    empirical estimate of generalization error.
    It captures the combined effects of bias,
    variance, and irreducible noise without
    measuring them individually.

    \item \textbf{Regularization:}
    By restricting model flexibility,
    regularization can reduce variance and
    prevent overfitting, potentially at the
    cost of increased bias.
\end{enumerate}

\subsubsection{Modern Overparameterized Models}

Classical intuition suggests that increasing model
complexity eventually increases variance and
test error.

However, modern overparameterized models can
achieve near-zero training error while still
generalizing well.

This behavior, sometimes associated with
\textbf{double descent}, challenges the simple
classical bias--variance tradeoff picture,
although the decomposition itself remains valid
under its assumptions.

\section{Neural Networks 1}

\textbf{Continuing with the Housing Price Prediction Problem}

\begin{figure}[H]
    \centering
    \includegraphics[width=0.8\linewidth]{Image/nonlinearHousingPricePrediction.png}
    \caption{Housing Price Prediction Dataset}
    \label{fig:housing_price_prediction}
\end{figure}

Our goal is to find a hypothesis function $h_\theta(x)$ that accurately captures the underlying relationship between the input features and housing prices. However, a linear model may fail to capture the nonlinear patterns in the dataset, resulting in underfitting. Therefore, we will explore nonlinear models to better fit the data.

Consider the function:
\[
    max\{t + wb, 0 \}  + c
\]

or

\[
    \mathrm{ReLU}(x) + c
\]

This is the best suit for the dataset.

Now, let's try it with a higher dimension. Suppose we have data with dimension $d$

\begin{description}
    \item[Given: ] Dataset
    
    $x \in \mathbb{R}^d$
    
    $w \in \mathbb{R}^d (\text{weight vector})$ 

    $b \in \mathbb{R}$ (\text{bias})

    \item[Defined Model: ] $h_\theta(x) = \mathrm{ReLU}(w^Tx + b) \in \mathbb{R}$ (single neuron) 
    
\end{description}

The main idea of a neural network is to learn useful nonlinear representations of the input data and combine them to make predictions.

Consider a neural network with three hidden neurons:

\[
\begin{aligned}
a_1 &= \operatorname{ReLU}(w_1^T x + b_1) \\
a_2 &= \operatorname{ReLU}(w_2^T x + b_2) \\
a_3 &= \operatorname{ReLU}(w_3^T x + b_3)
\end{aligned}
\]

Each neuron learns a nonlinear feature representation of the input. These features are then combined linearly at the output layer:

\[
\boxed{
h_\theta(x) = v_1a_1 + v_2a_2 + v_3a_3 + c
}
\]

Although the output is a linear combination of the learned features, the overall model remains nonlinear with respect to the input $x$.

\subsubsection{Why Neural Networks?}

Unlike linear regression, which directly combines input features, neural networks automatically learn useful feature representations through training.

\begin{itemize}
    \item \textbf{Feature Learning:} Hidden neurons learn useful representations of the input data.
    
    \item \textbf{Nonlinearity:} Activation functions allow the network to capture complex nonlinear relationships.
    
    \item \textbf{Feature Combination:} The output layer combines learned features to produce predictions.
    
    \item \textbf{Training:} Weights and biases are optimized using backpropagation and gradient-based optimization.
\end{itemize}

\textbf{Note:} A neuron does not necessarily represent a single interpretable feature. Multiple neurons may learn overlapping or distributed representations.

For deeper neural networks, multiple layers progressively transform and combine learned features:

\[
\boxed{
\text{Input} \rightarrow
\text{Feature Learning} \rightarrow
\text{Feature Combination} \rightarrow
\text{Prediction}
}
\]


\subsection{Back Propagation}

\[
    max(t, 0) = \text{relu} (t)
\]

\begin{description}
    \item[SGD: ]  

    \[
        J(\theta) = \frac{1}{n} \sum_{i = 1}^n J^{(i)}(\theta)
    \]

 And a loss function
    \[
        l_{cross\_entry} (h_\theta(x^{(i)}), y^{(i)})
    \]

    \[
        \theta = \theta_0 - \eta \nabla J^{(i_k)}(\theta)
    \]

 The main idea of back propagation is focusing on how we calculate $\nabla J^{(i_k)}(\theta)$ efficiently
\end{description}


\begin{description}
    \item[Differentiable Circuit / Network:]
    A differentiable circuit is a directed acyclic graph (DAG)
    consisting of:
    \begin{itemize}
        \item \textbf{Input nodes:} Variables $x_1, x_2, \ldots, x_l$.
        \item \textbf{Operation nodes:} Differentiable operations
        such as addition, multiplication, and activation functions.
        \item \textbf{Edges:} Connections that represent the flow
        of values between operations.
        \item \textbf{Output node:} The final computed value $f(x)$.
    \end{itemize}

    \textbf{Circuit vs. Expression:}
    The key difference is that a circuit can reuse intermediate
    results, while an expression tree may need to recompute
    or repeat the same subexpression multiple times.

    Therefore, a circuit can represent certain computations
    more efficiently than an expression tree.
\end{description}


\begin{description}
    \item[Theorem: ] (informal)

        Suppose a differentiable circuit of size $N$ (size means how many operations you have) computes a real-valued function  $f: \mathbb{R}^l \rightarrow \mathbb{R}$. Then, the gradient $\nabla f(x) \in \mathbb{R}^l$ can also be computed in time $O(N)$ by a circuit of size $O(N)$
\end{description}


\begin{description}
    \item[Backpropagation:]
    Backpropagation computes gradients efficiently by
    applying the \textbf{Chain Rule} backward through
    a computational graph.

    \textbf{Matrix Form of the Chain Rule:}

    Suppose $J$ depends on intermediate variables
    $a_1,\ldots,a_m$, which depend on $z_1,\ldots,z_n$.

    \[
    \underbrace{
    \begin{bmatrix}
        \frac{\partial J}{\partial z_1} &
        \cdots &
        \frac{\partial J}{\partial z_n}
    \end{bmatrix}
    }_{\text{Gradient w.r.t. } \mathbf z}
    =
    \underbrace{
    \begin{bmatrix}
        \frac{\partial J}{\partial a_1} &
        \cdots &
        \frac{\partial J}{\partial a_m}
    \end{bmatrix}
    }_{\text{Upstream Gradient}}
    \underbrace{
    \begin{bmatrix}
        \frac{\partial a_1}{\partial z_1} & \cdots &
        \frac{\partial a_1}{\partial z_n}\\
        \vdots & \ddots & \vdots\\
        \frac{\partial a_m}{\partial z_1} & \cdots &
        \frac{\partial a_m}{\partial z_n}
    \end{bmatrix}
    }_{\text{Local Jacobian}}
    \]

    In compact notation (row-vector convention):
    \[
        \boxed{
        \frac{\partial J}{\partial \mathbf z}
        =
        \frac{\partial J}{\partial \mathbf a}
        \frac{\partial \mathbf a}{\partial \mathbf z}
        }
    \]

    \textbf{Key Idea:}
    The gradient of each variable is obtained by
    multiplying the upstream gradient by the local
    Jacobian. Backpropagation applies this operation
    repeatedly in reverse order, reusing intermediate
    gradients instead of computing each derivative
    independently.
\end{description}

\chapter{Deep Learning for Computer Vision}


\section{Image Classification with Linear Classifiers}

\subsection{Image Classification: A core task in Computer Vision}

\begin{figure}[H]
    \centering
    \includegraphics[width=0.8\linewidth]{Image/WhatComputerSee.png}
    \label{fig:placeholder}
\end{figure}

Image are often define as matrices of data, each of pixels value are between $[0, 255]$

\subsection{Nearest Neighbor Classifiers}

\begin{figure}[H]
    \centering
    \includegraphics[width=0.8\linewidth]{Image/FirstClassifer.png}
    \label{fig:placeholder}
\end{figure}

\begin{description}
    \item[Distance Metrics to compare images: ] L1 distances: 
    \[
        d_1(I_1, I_2) = \sum_{p} |I_1^p - I_2^p|
    \]
\end{description}
\begin{figure}[H]
    \centering
    \includegraphics[width=0.8\linewidth]{Image/Metrics.png}
    \label{fig:placeholder}
\end{figure}

The implementation for the prediction:
\begin{figure}[H]
    \centering
    \includegraphics[width=0.8\linewidth]{Image/Implementation.png}
    \caption{Caption}
    \label{fig:placeholder}
\end{figure}

\subsection{K-Nearest Neighbors}

Instead of copying the label from the nearest neighbor, take the majority vote from $K$ closest points

\begin{description}
    \item[Distance Metrics: ] L2 distances:
    \[
        d_2(I_1, I_2) = \sqrt{\sum_p (I^p_1 - I^p_2)^2} 
    \]
    \item[Hyperparameters: ] What

    is the best value of $k$ to use?

    is the best distance to use?
    
\end{description}

\subsection{Linear Classifier}

A Linear Classifier is one of the fundamental and most important models in Machine Learning and Computer Vision. It uses a linear function to map input data to class scores.

The general formula of a linear classifier is defined as:
\begin{equation}
    f(x, W) = Wx + b
\end{equation}

Where the components are:
\begin{itemize}
    \item \textbf{$x$ (Input):} An input vector representing the data (e.g., an image).
    \item \textbf{$W$ (Weights / Parameters):} The weight matrix that determines the importance of input features for each class.
    \item \textbf{$b$ (Bias):} A bias vector that shifts the decision boundary independently of the input data $x$.
    \item \textbf{$f(x, W)$ (Class Scores):} A vector containing predicted scores for each class. A higher score indicates higher confidence that the data belongs to that class.
\end{itemize}

To deeply understand the nature of a linear classifier, we examine it through 3 viewpoints:

\begin{enumerate}
    \item \textbf{Algebraic Viewpoint:}
    The core operation is a matrix multiplication between the weights $W$ and the input vector $x$, followed by adding the bias $b$ to compute independent scores for each class.
    
    \item \textbf{Visual Viewpoint:}
    Each row of the weight matrix $W$ acts as a \textbf{"template"} representing a class. The score is effectively a measure of similarity between the input image $x$ and that class's template.
    
    \item \textbf{Geometric Viewpoint:}
    From a geometric perspective, a linear classifier partitions the feature space into regions using \textbf{hyperplanes} to separate different data classes.
\end{enumerate}

\section{Regularization and Optimization}

\subsection{Regularization}

We have the following loss function:
\[
    L(W) = \frac{1}{N} \sum_{i = 1}^{N}L_i(f(x_i, W), y_i)
\]

\textbf{Data loss:} Model predictions should match training data

However, the loss function included another component:
\[
    L(W) = \frac{1}{N} \sum_{i = 1}^{N}L_i(f(x_i, W), y_i) + \lambda R(W)
\]

\textbf{Regularization: } Prevent the model from doing \texttt{too well} on data.

\begin{figure}[H]
    \centering
    \includegraphics[width=0.8\linewidth]{Image/REgulazationIntuition.png}
    \caption{Caption}
    \label{fig:placeholder}
\end{figure}

\begin{description}
    \item[L2 Regularization (Ridge):] 
    Penalty term formula[cite: 2]:
    \begin{equation}
        R(W) = \sum_k \sum_l W_{k,l}^2
    \end{equation}
    \textit{Characteristics:} Computes the sum of squared elements in the weight matrix $W$[cite: 2]. Heavily penalizes large weights, driving them to distribute evenly and approach zero without ever becoming exactly zero, resulting in a smoother model.

    \item[L1 Regularization (Lasso):] 
    Penalty term formula[cite: 2]:
    \begin{equation}
        R(W) = \sum_k \sum_l |W_{k,l}|
    \end{equation}
    \textit{Characteristics:} Computes the sum of absolute values of elements in the weight matrix $W$[cite: 2]. Tends to drive less important weights to \textbf{exactly zero}, facilitating feature selection and producing a sparse model.

    \item[Elastic Net (L1 + L2):] 
    Combined formula[cite: 2]:
    \begin{equation}
        R(W) = \sum_k \sum_l (\beta W_{k,l}^2 + |W_{k,l}|)
    \end{equation}
    \textit{Characteristics:} Combines the advantages of both L1 and L2 via a coefficient $\beta$[cite: 2], offering strong feature selection capabilities while maintaining stability when features are highly correlated.
\end{description}

\subsection{Optimization}

First of all, we have 3 main variants of optimization methods:

\begin{enumerate}
    \item \textbf{Standard SGD (Stochastic Gradient Descent):} Updates weights based on individual data samples.
    \item \textbf{Batch Gradient Descent:} Computes gradients using the entire dataset in each iteration.
    \begin{quote}
        \textit{Note:} Since Standard SGD and Batch SGD were already covered in Machine Learning, we will focus a bit more on the next improvement.
    \end{quote}
    \item \textbf{SGD + Momentum:} An improved method that accelerates convergence and helps bypass local minima or saddle points by integrating inertia from previous steps.
\end{enumerate}

\subsubsection{Focusing on: SGD + Momentum}
The mathematical formula describing the weight update process in SGD + Momentum:
\begin{align}
    v_{t} &= \gamma v_{t-1} + \eta \nabla L(W_t) \\
    W_{t+1} &= W_t - v_t
\end{align}

Where the components are:
\begin{description}
    \item[$v_t$:] The velocity vector at step $t$, which retains the movement trend from previous steps.
    \item[$\gamma$:] The momentum coefficient (typically set to values like $0.9$ or $0.99$).
    \item[$\eta$:] The learning rate.
    \item[$\nabla L(W_t)$:] The gradient of the loss function at the current weights $W_t$.
\end{description}

\begin{figure}[H]
    \centering
    \includegraphics[width=0.8\linewidth]{Image/sgdmomentum.png}
    \label{fig:placeholder}
\end{figure}

\subsubsection{More complex optimizers: RMSProp}

Currently, I find it is not necessary to learn this method. However, it would be useful in the future when I try to optimize my model, and so on. So, I will note here but not completely understand this expression. 

\begin{figure}[H]
    \centering
    \includegraphics[width=0.8\linewidth]{Image/RMSProp.png}
    \label{fig:placeholder}
\end{figure}

\subsection{Optimizers: Adam(almost)}

This is the most popular optimizers in deep learning field.

\begin{figure}[H]
    \centering
    \includegraphics[width=0.9\linewidth]{Image/Adam.png}
    \caption{Implementation of the Optimzier}
    \label{fig:placeholder}
\end{figure}

\subsubsection{Learning Rate}

\begin{figure}[H]
    \centering
    \includegraphics[width=0.9\linewidth]{Image/lrdecay.png}
    \label{fig:placeholder}
\end{figure}

\begin{description}
    \item[Step: ] Reduce learning rate at a few fixed points. E.g. for ResNets, multiply LR by $0.1$ after epochs 30, 60, and 90.

    \begin{itemize}
        \item \textbf{Cosine} $\alpha_t = \frac{1}{2}\alpha_0 (1 + \cos(t\pi/T))$

        \item \textbf{Linear} $\alpha_t = \alpha_0(1 - t/T)$

        \item \textbf{Inverse sqrt} $\alpha_t = \alpha/\sqrt t$
    \end{itemize}

    In which: 
    \[
        \alpha_0: \text{ Initial Learning Rate} \\
        \alpha_t: \text{ Learning rate at epoch t} \\
        T: \text{ Total number of epochs}
    \]
\end{description}

\section{Better Visual in Backpropagation and Neural Network}

\begin{figure}[H]
    \centering
    \includegraphics[width=0.8\linewidth]{Image/vectorderivate.png}
    \label{fig:placeholder}
\end{figure}

\subsection{Backpropagation with Vector}

\begin{figure}[H]
    \centering
    \includegraphics[width=0.8\linewidth]{Image/backpropwvector.png}
    \label{fig:placeholder}
\end{figure}

\section{Image Classification with CNNs}



\end{document}
